\documentclass[11pt, twocolumn]{article}
\usepackage{graphicx}
\usepackage{multicol}
\usepackage{tabularx}
\usepackage[margin=1in]{geometry}
\usepackage{url}
\usepackage{listings}
%\usepackage[style=ieee]{biblatex}

\renewcommand{\abstractname}{Resumen}
\renewcommand{\refname}{Referencias}

\begin{document}

\begin{titlepage}
	\centering
	\LARGE
	\textbf{INSTITUTO TECNOLÓGICO Y DE ESTUDIOS SUPERIORES DE OCCIDENTE}
	
	\large
	\textbf{Departamento de Electrónica, Sistemas e Informática}\\
	\textbf{Maestría en Sistemas Computacionales}\\
	\vspace{1.0cm}
	\textbf{Procesamiento digital de señales}
	\textbf{Otoño 2022}
	\includegraphics[scale=1.2]{images/iteso_logo}

	\LARGE
	\textbf{Tarea 0: Resumen de artículos sobre Procesamiento de Señales}
	\vfill
	
	\textsf{741645, erich.cordoba@iteso.mx}\\
	\textsc{Tlaquepaque, Jalisco. Agosto 2022}
\end{titlepage}


\begin{abstract}
El presente trabajo resume el contenido de tres artículos sobre el Procesamiento Digital de Señales (PDS). Los artículos analizados se han seleccionado de distintos momentos en la historia reciente del PDS, destacando la visión y estimación puntual sobre las problemáticas y tendencias esperadas en el desarrollo de esta rama de investigación.
\end{abstract}

\section{Introducción}
La evolución del PDS ha estado ligada diréctamente al desarrollo tecnológico de los dispositivos que la hacen posible. En los años previos a la década de 1950, el PDS se realizaba únicamente con circuitos analógicos. Posteriormente, la inclusión de circuitos digitales benefició el modelado del PDS, sin embargo el desempeño de estos circuitos limitó durante un tiempo la implementación de aplicaciones de PDS. Es en la últimas dos décadas donde se ha presentado el mayor incremento en la investigación y desarrollo de PDS. Cada vez se tiene acceso a nuevos dispositivos con mejoras en prestaciones de procesamiento y consumo de energía, microcontroladores optimizados para realizar operaciones de multiplicación y suma de una forma eficiente. Así mismo, es posible diseñar circuitos a la medida mediante FPGAs y ASICs. Incluso se puede considerar el uso de unidades de procesamiento de gráficos para alguna aplicación de PDS.\\

Es de espera que, durante estos últimos años de un incremento considerable en el desarrollo de PDS, nuevas problemáticas estén surgiendo demandando técnicas y soluciones innovadoras. Haykin \cite{haykin}, presentó en 2001 un artículo describiendo el estado y retos futuros en el PDS. En particular, Haykin destacó como problema principal la separación entre las matemáticas y física en el diseño de soluciones de PDS. Años más tarde, Deng \cite{golden_age}, visita las aplicaciones principales de PDS en 2009, mencionando la necesidad de crear grupos interdisciplinarios que ayuden a comprender mejor las señales que se están interpretando. Finalmente en 2020, Brower \cite{dps_dead} hace la pregunta si los procesadores digitales de señales han muerto, respondiendo a detalle con las aplicaciones actuales y la tendencia marcada a incorporar técnicas de inteligencia artificial.\\

Las siguientes secciones resumen cada uno de estos tres artículos.

\section{Procesamiento de señales: cuando las matemáticas y la física se encuentran}
Descripción del primer artíclo

\section{Abrazando una nueva era dorada en el procesamiento de señales}
Hablar sobre el segundo artículo

\section{¿Están los PDSs muertos?}
Hablar sobre el tercer artículo


\section{Referencias}

\bibliographystyle{ieeetr}
\bibliography{referencias}

\end{document}